\section*{Bloque D. Ampliación: Teoría y Aplicaciones Prácticas (70\% / 30\%)}

\noindent\textit{Preguntas 91 a 180. Ampliación del banco oficial con 63 preguntas teóricas y conceptuales (91 a 153) y 27 aplicaciones cuantitativas prácticas (154 a 180).}

\subsection*{Parte I. Preguntas teóricas y conceptuales (91 a 153)}

\q{En un juego estático finito con información completa, la eliminación iterativa de estrategias estrictamente dominadas (EIEED) garantiza que:}
\begin{opts}
    \item El resultado final es independiente del orden de eliminación y contiene todos los Equilibrios de Nash del juego.
    \item El orden en que se eliminan las estrategias estrictamente dominadas puede alterar el conjunto final de perfiles sobrevivientes.
    \item Siempre conduce a un único perfil de estrategias puras si los pagos de todos los jugadores son estrictamente positivos.
    \item Elimina únicamente estrategias puras y nunca puede eliminar una estrategia dominada por una combinación de estrategias mixtas.
\end{opts}

\q{A diferencia de la eliminación iterativa de estrategias estrictamente dominadas, la eliminación iterativa de estrategias débilmente dominadas presenta la propiedad de que:}
\begin{opts}
    \item Siempre converge al mismo equilibrio único en estrategias puras sin importar qué jugador elimine primero sus acciones.
    \item El conjunto sobreviviente puede depender del orden de eliminación y descartar Equilibrios de Nash.
    \item Garantiza la preservación de todos los Equilibrios de Nash en estrategias puras y mixtas del juego original.
    \item Requiere únicamente racionalidad individual de primer orden y no necesita conocimiento común de la racionalidad.
\end{opts}

\q{En el juego de la Caza del Ciervo (Stag Hunt), existen dos Equilibrios de Nash en estrategias puras: (Ciervo, Ciervo) y (Liebre, Liebre). Respecto a ambos equilibrios, es correcto afirmar que:}
\begin{opts}
    \item El equilibrio (Liebre, Liebre) es Pareto-eficiente, mientras que (Ciervo, Ciervo) es dominante en estrategias estrictas.
    \item Ambos equilibrios son estrictamente Pareto-comparables y (Liebre, Liebre) domina en pagos al equilibrio cooperativo.
    \item El equilibrio (Ciervo, Ciervo) domina en pagos (Pareto-superior), mientras que (Liebre, Liebre) puede dominar en riesgo.
    \item No existe ningún Equilibrio de Nash en estrategias mixtas porque las estrategias puras forman equilibrios estrictos.
\end{opts}

\q{De acuerdo con el Teorema de Existencia de Nash (1950), todo juego no cooperativo en forma normal con un número finito de jugadores y un conjunto finito de estrategias puras:}
\begin{opts}
    \item Posee necesariamente al menos un Equilibrio de Nash en estrategias puras si la suma de pagos es constante.
    \item Tiene exactamente un Equilibrio de Nash en estrategias mixtas cuando no existe ninguna estrategia estrictamente dominante.
    \item Garantiza que todos sus Equilibrios de Nash son eficientes en el sentido de Pareto si los jugadores son racionales.
    \item Posee al menos un Equilibrio de Nash cuando se permite que los jugadores utilicen estrategias mixtas.
\end{opts}

\q{En un juego $2\times 2$ con un único Equilibrio de Nash en estrategias completamente mixtas, si aumenta el pago del Jugador 1 asociado a una de sus acciones puras (manteniendo constantes los pagos del Jugador 2):}
\begin{opts}
    \item Cambia la mezcla del Jugador 2 para restaurar la indiferencia del Jugador 1, sin alterar la mezcla de este.
    \item Aumenta directamente la probabilidad con la que el Jugador 1 elige esa acción, sin alterar la mezcla óptima del Jugador 2.
    \item Ambos jugadores incrementan la probabilidad de jugar sus primeras acciones puras en la misma proporción marginal.
    \item El juego deja inmediatamente de tener un Equilibrio de Nash en estrategias mixtas y converge a un equilibrio en puras.
\end{opts}

\q{En un juego en forma extensiva con información imperfecta, un subjuego propio debe comenzar necesariamente en un nodo de decisión que:}
\begin{opts}
    \item Corresponda exclusivamente al primer movimiento del juego realizado por la naturaleza o por el jugador líder.
    \item Pertenezca a un conjunto de información unitario y no divida ningún conjunto de información sucesor.
    \item Contenga al menos dos nodos de decisión indistinguibles para que el jugador deba formar creencias bayesianas.
    \item Conduzca de manera inmediata a nodos terminales sin permitir movimientos adicionales de otros jugadores.
\end{opts}

\q{En el juego secuencial de entrada al mercado, una empresa incumbente amenaza con desencadenar una guerra de precios destructiva si un entrante ingresa. El Equilibrio de Nash Perfecto en Subjuegos (ENPS) descarta el equilibrio donde el entrante no ingresa porque:}
\begin{opts}
    \item La amenaza del incumbente viola la restricción de participación del entrante en el nodo inicial del árbol de juego.
    \item El entrante carece de información perfecta sobre los costos fijos hundidos del incumbente antes de decidir su entrada.
    \item La guerra de precios no es una mejor respuesta del incumbente en el subjuego que se inicia tras producirse la entrada.
    \item En todo juego secuencial de dos etapas existe siempre un único Equilibrio de Nash en estrategias puras sin refinamientos.
\end{opts}

\q{En teoría de juegos dinámicos, el valor estratégico de un compromiso irreversible (como instalar capacidad excesiva o firmar un contrato con penalidad) radica en que:}
\begin{opts}
    \item Amplía el conjunto de opciones disponibles para el jugador en etapas posteriores, garantizando mayor flexibilidad operativa.
    \item Elimina la interdependencia estratégica entre rivales al transformar el oligopolio en un mercado perfectamente competitivo.
    \item Reduce el factor de descuento intertemporal de los competidores y elimina la necesidad de inducción hacia atrás.
    \item Altera los incentivos propios futuros, volviendo creíble una amenaza y modificando la respuesta del rival.
\end{opts}

\q{Si un Dilema del Prisionero con un único Equilibrio de Nash en el juego de etapa se repite un número finito y conocido de veces $T$, la aplicación de inducción hacia atrás implica que:}
\begin{opts}
    \item El único Equilibrio de Nash Perfecto en Subjuegos consiste en jugar la acción no cooperativa en todos y cada uno de los $T$ períodos.
    \item La cooperación puede sostenerse durante los primeros $T-1$ períodos siempre que el factor de descuento sea suficientemente cercano a $1$.
    \item Los jugadores cooperan en la última etapa $T$ para evitar el castigo de la estrategia gatillo en períodos futuros.
    \item Existen múltiples ENPS eficientes en el sentido de Pareto gracias al Teorema Popular (Folk Theorem) para juegos finitos.
\end{opts}

\q{En un juego repetido con horizonte indefinido donde en cada período existe una tasa de interés real $r > 0$ y una probabilidad constante $p \in (0,1)$ de que el juego continúe al siguiente período, el factor de descuento efectivo $\delta$ es igual a:}
\begin{opts}
    \item $\delta = p(1 + r)$, de modo que una mayor tasa de interés facilita la sostenibilidad de la colusión tácita.
    \item $\delta = \frac{p}{1 + r}$, reflejando tanto la impaciencia financiera como el riesgo de terminación del juego.
    \item $\delta = \frac{1 - p}{1 + r}$, disminuyendo cuando aumenta la probabilidad de que las empresas sigan interactuando.
    \item $\delta = \frac{1 + r}{p}$, siendo necesariamente mayor que la unidad para que converja el valor presente.
\end{opts}

\q{Respecto a la estrategia del gatillo implacable (Grim Trigger) frente a esquemas de castigo óptimo en dos fases (tipo 'palo y zanahoria' de Abreu) en juegos repetidos de Cournot:}
\begin{opts}
    \item Grim Trigger siempre impone el castigo más severo posible y logra el umbral mínimo absoluto del factor de descuento.
    \item Los castigos de dos fases prohíben retornar a la colusión tras una desviación, haciéndolos menos creíbles que Grim Trigger.
    \item Un castigo severo y temporal seguido de retorno a la cooperación puede sostener la colusión con un $\delta$ menor que Grim Trigger.
    \item Ningún esquema de castigo puede resultar perfecto en subjuegos si exige que la empresa castigadora sacrifique beneficios temporales.
\end{opts}

\q{En un juego bayesiano estático con información incompleta (transformación de Harsanyi), una estrategia de un jugador se define formalmente como:}
\begin{opts}
    \item Una única acción constante que el jugador elige sin condicionarla a la información privada que posee sobre su propio tipo.
    \item Un vector de pagos esperados que garantiza inmunidad frente a desviaciones colusivas de coaliciones de jugadores.
    \item Una regla de actualización de creencias que aplica el teorema de Bayes únicamente en nodos fuera de la senda de equilibrio.
    \item Un plan contingente que asigna una acción (o mezcla) a cada uno de los posibles tipos privados que ese jugador puede tener.
\end{opts}

\q{En el juego de la Guerra de Desgaste (War of Attrition) donde dos empresas deciden cuánto tiempo permanecer en un mercado que solo es rentable para una firma:}
\begin{opts}
    \item Existen dos equilibrios asimétricos en puras (una sale y otra queda) y uno simétrico en estrategias mixtas.
    \item La salida simultánea inmediata de ambas empresas constituye el único Equilibrio de Nash en estrategias puras del juego.
    \item Las empresas siempre se reparten el mercado de manera equitativa sin incurrir en pérdidas acumuladas por permanencia.
    \item Aumentar el costo por período de permanencia induce a ambas empresas a quedarse indefinidamente en el mercado.
\end{opts}

\q{El Principio de Desviación en una Etapa (One-Shot Deviation Principle) establece que, en un juego dinámico finito o infinito con pagos descontados continuos al infinito, un perfil de estrategias es ENPS si y solo si:}
\begin{opts}
    \item Ningún jugador puede mejorar su pago esperado desviándose de manera permanente durante todos los períodos restantes.
    \item Ningún jugador gana desviándose en un único período (o conjunto de información) y retomando luego el plan.
    \item Todos los jugadores utilizan estrategias estacionarias independientes de la historia previa del juego.
    \item La suma de los pagos de todos los jugadores alcanza el máximo bienestar social en cada subjuego del árbol.
\end{opts}

\q{En una subasta de valor común (por ejemplo, derechos de exploración o concesiones de infraestructura con costos inciertos compartidos), la 'maldición del ganador' (winner's curse) ocurre porque:}
\begin{opts}
    \item Los postores con mayor aversión al riesgo pujan sistemáticamente por encima de su valoración esperada incondicional.
    \item El ganador paga necesariamente el segundo precio más alto, lo cual supera siempre el valor intrínseco del activo subastado.
    \item Quien gana suele tener la señal más optimista, revelando malas noticias sobre el valor real si no sombreó su puja.
    \item El subastador revela información privada después de adjudicar el contrato, reduciendo el valor de mercado del proyecto.
\end{opts}

\q{Según el Teorema de Equivalencia de Ingresos (Vickrey-Myerson), cuatro formatos estándar de subasta (primer precio, segundo precio, inglesa y holandesa) generan el mismo ingreso esperado para el vendedor cuando:}
\begin{opts}
    \item Los postores son fuertemente aversos al riesgo y poseen valoraciones comunes correlacionadas positivamente en equilibrio.
    \item Las valoraciones son conocidas públicamente antes de la licitación y los costos de participación son asimétricos.
    \item Existe colusión tácita entre los postores y el subastador fija un precio de reserva estrictamente igual a cero.
    \item Los postores son neutrales al riesgo, tienen valores privados independientes simétricos y gana el de mayor valor.
\end{opts}

\q{En la clasificación de Bulow, Geanakoplos y Klemperer (1985), las variables estratégicas $(x_1, x_2)$ de dos oligopolistas son complementos estratégicos cuando:}
\begin{opts}
    \item $\frac{\partial^2 \pi_i}{\partial x_i \partial x_j} > 0$, dando funciones de reacción crecientes como en Bertrand diferenciado.
    \item $\frac{\partial \pi_i}{\partial x_j} > 0$, sin que importe el signo de la derivada cruzada sobre la función de mejor respuesta.
    \item $\frac{\partial^2 \pi_i}{\partial x_i \partial x_j} < 0$, lo que produce funciones de reacción con pendiente negativa como en Cournot lineal.
    \item $\frac{\partial^2 \pi_i}{\partial x_i^2} > 0$, garantizando que las curvas de isobeneficio sean estrictamente convexas respecto al origen.
\end{opts}

\q{En un duopolio de Cournot con demanda lineal y costos marginales constantes $(c_1, c_2)$, una innovación tecnológica que reduce únicamente el costo marginal de la empresa 1 ($\downarrow c_1$) provoca en el nuevo equilibrio:}
\begin{opts}
    \item Un aumento proporcional tanto de $q_1$ como de $q_2$, manteniendo inalteradas las cuotas de mercado de ambas empresas.
    \item Un aumento de $q_1$, una reducción de $q_2$, un aumento de la producción total $Q$ y una caída del precio de mercado $P$.
    \item Un aumento de $q_1$ compensado exactamente por la caída de $q_2$, dejando sin cambios la cantidad total $Q$ y el precio $P$.
    \item Una reducción del precio $P$ en exactamente la misma magnitud que la caída de $c_1$, eliminando el margen de la empresa 2.
\end{opts}

\q{En un oligopolio de Cournot con $N$ empresas simétricas, demanda lineal y costo marginal constante $c$, conforme el número de empresas $N$ crece hacia el infinito:}
\begin{opts}
    \item El precio converge al precio de monopolio porque las empresas coordinan implícitamente sus cuotas de mercado.
    \item El beneficio individual por empresa disminuye, pero la suma de beneficios de la industria permanece constante en el nivel monopólico.
    \item La cantidad total converge al nivel competitivo $Q_C$, el precio converge a $c$ y la pérdida de eficiencia social tiende a cero.
    \item La pendiente de las funciones de reacción se vuelve positiva, transformando la competencia en cantidades en Bertrand diferenciado.
\end{opts}

\q{En un oligopolio de Cournot con $N$ empresas que producen un bien homogéneo con costos marginales $c_i$ y cuotas de mercado $s_i = q_i/Q$, el Índice de Lerner promedio ponderado de la industria ($\bar{L} = \sum s_i L_i$) satisface:}
\begin{opts}
    \item $\bar{L} = \frac{CR_2}{N \cdot |\varepsilon_d|}$, dependiendo exclusivamente de las dos empresas con mayor participación de mercado.
    \item $\bar{L} = HHI \cdot |\varepsilon_d|$, indicando que una demanda más elástica eleva el margen precio-costo de las empresas.
    \item $\bar{L} = \frac{1 - HHI}{|\varepsilon_d|}$, alcanzando su valor máximo cuando el mercado está completamente atomizado.
    \item $\bar{L} = \frac{HHI}{|\varepsilon_d|}$, siendo $HHI = \sum s_i^2$ y $|\varepsilon_d|$ la elasticidad de demanda.
\end{opts}

\q{En un oligopolio de Cournot simétrico con $N$ empresas, demanda lineal y costo marginal constante $c$, la imposición de un impuesto específico de $t$ soles por unidad producida genera un incremento en el precio de equilibrio igual a:}
\begin{opts}
    \item $\Delta P = \frac{N}{N+1} t$, de modo que la traslación al precio es parcial y crece hacia $t$ conforme aumenta $N$.
    \item $\Delta P = \frac{t}{N+1}$, disminuyendo la traslación hacia los consumidores a medida que aumenta la competencia.
    \item $\Delta P = t$, trasladándose el 100\% del impuesto a los consumidores independientemente del número de competidores.
    \item $\Delta P = \frac{N+1}{N} t$, produciendo siempre una sobretraslación (overshooting) superior al monto del impuesto.
\end{opts}

\q{En un duopolio de Bertrand con bienes homogéneos, sin restricciones de capacidad y costos marginales constantes asimétricos $c_1 < c_2$ (donde $c_2$ es menor que el precio de monopolio de la firma 1), el equilibrio se caracteriza porque:}
\begin{opts}
    \item Ambas empresas fijan un precio igual al promedio de sus costos marginales $\frac{c_1 + c_2}{2}$ y se dividen la demanda.
    \item La empresa más eficiente fija un precio límite $P^* = c_2$, abastece todo el mercado y obtiene beneficios positivos.
    \item La empresa menos eficiente fija $P = c_1$ para expulsar a su rival y luego recupera sus pérdidas mediante colusión tácita.
    \item No existe equilibrio ni siquiera aproximado porque la empresa 1 siempre preferirá fijar su precio de monopolio irrestricto.
\end{opts}

\q{En el modelo de Bertrand-Edgeworth, cuando dos empresas que venden un bien homogéneo compiten en precios pero enfrentan restricciones de capacidad intermedias (ninguna puede abastecer sola toda la demanda a $P = c$):}
\begin{opts}
    \item El precio de equilibrio cae por debajo del costo marginal $c$ para liquidar la capacidad ociosa en el corto plazo.
    \item Las empresas necesariamente alcanzan el resultado de monopolio conjunto sin necesidad de repetir el juego en el tiempo.
    \item El perfil $P_1 = P_2 = c$ deja de ser Equilibrio de Nash en puras y surgen equilibrios en estrategias mixtas.
    \item La empresa con menor capacidad instalada captura el 100\% del mercado fijando un precio igual al costo medio mínimo.
\end{opts}

\q{El resultado fundamental del juego en dos etapas de Kreps y Scheinkman (1983) —donde las empresas eligen primero su capacidad de producción y luego compiten en precios a la Bertrand— demuestra que:}
\begin{opts}
    \item La competencia en precios de la segunda etapa elimina cualquier margen, replicando siempre el equilibrio competitivo $P = c$.
    \item La elección de capacidad en la primera etapa siempre conduce al nivel de producción de un cártel perfectamente colusivo.
    \item Las empresas nunca instalan capacidad positiva en la primera etapa debido al riesgo de guerra de precios ex post.
    \item Bajo una regla de racionamiento eficiente, el equilibrio perfecto en subjuegos reproduce las cantidades y precios de Cournot.
\end{opts}

\q{En un duopolio de Bertrand con bienes diferenciados y demandas lineales $q_i = \alpha - \beta p_i + \gamma p_j$ ($0 < \gamma < \beta$), una reducción en el parámetro de sustituibilidad $\gamma$ (mayor diferenciación) provoca:}
\begin{opts}
    \item Un mayor poder de mercado para cada firma, con precios más elevados respecto al costo marginal y mayores márgenes.
    \item Un aumento en la pendiente de las funciones de reacción y una caída de los beneficios de equilibrio de ambas empresas.
    \item La desaparición de la demanda residual de la firma con mayor costo marginal y la salida automática del mercado.
    \item Una disminución del precio de equilibrio y una convergencia más rápida hacia la paradoja de Bertrand ($P = c$).
\end{opts}

\q{Al comparar el liderazgo secuencial (Stackelberg) en cantidades bajo sustitutos estratégicos frente al liderazgo secuencial en precios bajo complementos estratégicos (bienes diferenciados), se cumple que:}
\begin{opts}
    \item En ambos casos el líder obtiene estrictamente mayores beneficios que el seguidor gracias a la ventaja del primer movedor.
    \item En cantidades hay ventaja del primer movedor, pero en precios diferenciados suele haber ventaja del segundo.
    \item En precios diferenciados el líder fija un precio inferior al de Bertrand simultáneo para expulsar al seguidor del mercado.
    \item En ambos modelos el equilibrio secuencial produce exactamente los mismos precios y beneficios que el equilibrio simultáneo.
\end{opts}

\q{En teoría del oligopolio, el denominado 'desequilibrio de Bowley' (o guerra de Stackelberg) surge cuando en un duopolio en cantidades:}
\begin{opts}
    \item Ambas empresas actúan simultáneamente como seguidoras de Cournot, produciendo por debajo del nivel competitivo.
    \item El seguidor desconoce la función de demanda del mercado y adopta una regla de precios igual al costo medio.
    \item Ambas firmas intentan actuar como líderes de Stackelberg, generando sobreproducción y menores beneficios.
    \item Una empresa elige cantidad mientras la otra elige precio simultáneamente en un mercado de bien homogéneo.
\end{opts}

\q{Un cártel que busca maximizar los beneficios conjuntos de dos empresas con costos marginales crecientes y diferentes ($CMg_1(q_1)$ y $CMg_2(q_2)$) debe asignar las cuotas de producción de modo que:}
\begin{opts}
    \item La empresa de mayor costo marginal produzca una cuota mayor para compensar su menor rentabilidad operativa.
    \item Ambas empresas produzcan exactamente la mitad de la cantidad monopólica independientemente de sus estructuras de costos.
    \item Cada empresa iguale su costo medio mínimo al precio de monopolio para eliminar el incentivo a desviarse del acuerdo.
    \item $CMg_1(q_1) = CMg_2(q_2) = IMg(q_1 + q_2)$, igualando costos marginales entre plantas con el ingreso marginal.
\end{opts}

\q{En modelos de colusión tácita en juegos repetidos, ¿qué combinación de factores estructurales facilita (reduce el umbral $\delta^*$) la sostenibilidad de un acuerdo colusivo?}
\begin{opts}
    \item Menor número de empresas, interacción frecuente, alta transparencia de precios y simetría en costos y capacidades.
    \item Mayor número de empresas, pedidos grandes e infrecuentes y precios negociados en secreto entre proveedor y cliente.
    \item Barreras de entrada nulas, innovación tecnológica disruptiva continua y elevada heterogeneidad en los productos.
    \item Elevada tasa de interés real, alta probabilidad de quiebra sectorial y fuerte poder de negociación de los compradores.
\end{opts}

\q{En un oligopolio donde las empresas tienen costos marginales asimétricos, la colusión tácita suele resultar más difícil de sostener que en un mercado simétrico principalmente porque:}
\begin{opts}
    \item La empresa de mayor costo siempre puede capturar el 100\% del mercado desviándose unilateralmente a un precio igual a $c_1$.
    \item La firma de menor costo tiene mayor incentivo a desviarse y prefiere un precio colusivo menor que su rival.
    \item Las leyes antimonopolio prohíben únicamente los acuerdos horizontales entre empresas con distintos costos marginales.
    \item Los beneficios totales del cártel son menores que la suma de beneficios de Cournot cuando existen asimetrías de costos.
\end{opts}

\q{Según el análisis de Bernheim y Whinston (1990) sobre contacto multimercado (empresas que compiten simultáneamente en varios mercados geográficos o de productos), dicho contacto puede facilitar la colusión cuando:}
\begin{opts}
    \item Todos los mercados son completamente idénticos en demanda, costos y número de competidores con rendimientos constantes.
    \item Las empresas carecen de capacidad para observar las desviaciones de sus rivales en otros mercados geográficos.
    \item Existen asimetrías que permiten transferir holgura colusiva de un mercado estable hacia otro más difícil.
    \item El regulador obliga a fijar un precio único nacional igual al costo marginal de largo plazo en todos los mercados.
\end{opts}

\q{Aunque a primera vista parecen beneficiosas para el consumidor, las cláusulas de 'igualación de precios' (meet-the-competition) adoptadas por oligopolistas pueden actuar como mecanismos facilitadores de colusión porque:}
\begin{opts}
    \item Obligan legalmente a las empresas a fijar precios iguales al costo marginal de la empresa más eficiente del sector.
    \item Incrementan los costos de búsqueda de los consumidores y eliminan la transparencia de precios en el mercado minorista.
    \item Impiden que las empresas ofrezcan descuentos selectivos cuando la demanda agregada experimenta una recesión.
    \item Neutralizan la ganancia de recortar precios, pues la respuesta automática impide al rival robar clientes.
\end{opts}

\q{En el modelo de colusión con fluctuaciones observables de demanda de Rotemberg y Saloner (1986), los precios colusivos tienden a comportarse de manera contracíclica (bajando en auges) porque:}
\begin{opts}
    \item En auges la ganancia de desvío es alta frente al castigo futuro promedio, obligando a bajar el precio colusivo.
    \item Durante las recesiones las empresas tienen incentivos irresistibles a iniciar guerras de precios para liquidar inventarios.
    \item En períodos de auge el factor de descuento cae automáticamente a cero debido a la entrada inmediata de nuevos competidores.
    \item Las empresas confunden una caída inesperada de la demanda agregada con una desviación secreta de sus rivales.
\end{opts}

\q{En la política de defensa de la competencia (como el Programa de Clemencia de INDECOPI), conceder exoneración total de multa al primer miembro de un cártel que confiese y aporte pruebas busca principalmente:}
\begin{opts}
    \item Aumentar la recaudación fiscal de la agencia de competencia mediante multas más elevadas a las empresas no coludidas.
    \item Inducir un Dilema del Prisionero entre miembros del cártel, elevando la desconfianza y rompiendo la colusión.
    \item Permitir que las empresas coludidas se fusionen legalmente para aprovechar economías de escala sin sanción administrativa.
    \item Sustituir la regulación tarifaria de monopolios naturales por acuerdos voluntarios de autorregulación entre empresas.
\end{opts}

\q{Según la regla de Areeda-Turner en derecho de la competencia, un indicio clásico de fijación de precios predatorios por parte de una empresa dominante se configura cuando esta fija un precio:}
\begin{opts}
    \item Inferior al precio de monopolio de corto plazo, aunque superior a su costo medio total de largo plazo.
    \item Superior al costo marginal pero inferior al precio fijado por las empresas más pequeñas de la franja competitiva.
    \item Inferior al costo variable medio (aproximación del $CMg$ de corto plazo) para expulsar competidores.
    \item Igual al costo medio total con el propósito de disuadir la entrada en un mercado perfectamente desafiable.
\end{opts}

\q{En el modelo secuencial de disuasión de entrada de Dixit (1980), una empresa incumbente puede impedir de forma creíble el ingreso de un entrante mediante:}
\begin{opts}
    \item La promesa verbal de reducir su precio a cero si el entrante decide construir una planta en el segundo período.
    \item La contratación de deuda de corto plazo sin colateral para convencer al entrante de que evitará competir agresivamente.
    \item La fijación de un precio elevado antes de la entrada para señalar falsamente que el mercado tiene baja rentabilidad.
    \item La inversión irreversible en capacidad de producción en la etapa 1, que reduce su costo marginal ex post en la etapa 2.
\end{opts}

\q{En la taxonomía estratégica de Fudenberg y Tirole (1984) para mercados con acomodación de entrada, una estrategia de 'perro faldero' (Puppy Dog Ploy) es óptima cuando:}
\begin{opts}
    \item La inversión vuelve agresiva a la firma y los precios son complementos estratégicos, por lo que subinvierte.
    \item La inversión hace más agresivo al incumbente y las variables son sustitutos estratégicos, por lo que conviene sobreinvertir.
    \item El incumbente busca disuadir la entrada a toda costa reduciendo drásticamente sus precios antes del ingreso del rival.
    \item Las empresas compiten en un mercado perfectamente competitivo sin costos hundidos ni interdependencia estratégica.
\end{opts}

\q{La paradoja de las fusiones horizontales en Cournot (Salant, Switzer y Reynolds, 1983) demuestra que, en un mercado con $N \ge 3$ firmas simétricas, demanda lineal y costo marginal constante sin sinergias, la fusión de dos empresas:}
\begin{opts}
    \item Siempre eleva los beneficios conjuntos de las empresas fusionadas porque reduce la competencia y sube el precio de mercado.
    \item Suele reducir los beneficios de las fusionadas porque los rivales expanden su oferta al contraer ellas la suya.
    \item Reduce el precio de mercado y perjudica únicamente a las empresas competidoras que quedaron fuera de la fusión.
    \item Es rentable únicamente cuando el índice HHI previo a la fusión es inferior a $1000$ puntos.
\end{opts}

\q{A diferencia de Cournot, una fusión horizontal entre dos empresas que compiten en precios con productos diferenciados (complementos estratégicos) sin ganancias de eficiencia:}
\begin{opts}
    \item Resulta siempre no rentable para las fusionadas porque los competidores externos recortan agresivamente sus precios.
    \item Deja inalterados los precios de todos los productos del mercado al no modificar el número de plantas productivas.
    \item Suele ser rentable y ejerce presión al alza de precios (UPP), proporcional a la razón de desvío.
    \item Reduce necesariamente los beneficios de las empresas no participantes debido a la expansión de ventas de las fusionadas.
\end{opts}

\q{El problema de la doble marginalización ocurre en una estructura vertical de monopolios sucesivos (fabricante y distribuidor) con contratos lineales y se caracteriza porque:}
\begin{opts}
    \item Solo puede resolverse mediante la fijación de un precio mínimo de reventa (RPM) que impida descuentos minoristas.
    \item El precio final coincide con el costo marginal del fabricante, generando pérdidas económicas para el distribuidor minorista.
    \item La integración vertical eleva el precio final al consumidor al sumar los dos poderes de monopolio en una sola firma.
    \item El precio final supera al del monopolio integrado, reduciendo el excedente del consumidor y el beneficio total.
\end{opts}

\q{Desde la perspectiva económica de la Organización Industrial, ¿cuál es el principal argumento de eficiencia pro-competitiva a favor de la fijación de precios mínimos de reventa (RPM) o territorios exclusivos?}
\begin{opts}
    \item Evitan el free-riding entre minoristas, incentivando la provisión óptima de servicios pre-venta.
    \item Eliminan por completo el poder de mercado del fabricante al igualar el precio mayorista al costo marginal de producción.
    \item Permiten a los distribuidores locales formar un cártel minorista estable frente a un fabricante atomizado.
    \item Garantizan que el precio final al consumidor sea siempre inferior al costo marginal del distribuidor más eficiente.
\end{opts}

\q{En el modelo clásico de competencia monopolística de Chamberlin, la distinción entre la curva de demanda percibida por una firma ($d$) y la curva de demanda proporcional o de cuota de mercado ($D$) radica en que:}
\begin{opts}
    \item La curva $D$ es siempre perfectamente horizontal, mientras que $d$ tiene pendiente positiva debido a la diferenciación.
    \item La curva $d$ es más elástica al asumir precios rivales fijos, mientras $D$ asume ajustes simultáneos iguales.
    \item En equilibrio de largo plazo las curvas $d$ y $D$ nunca se intersectan porque los productos son sustitutos perfectos.
    \item La curva $d$ representa la demanda agregada de toda la industria y la curva $D$ refleja el costo marginal del productor.
\end{opts}

\q{En el equilibrio de largo plazo de competencia monopolística, la existencia de 'exceso de capacidad' ($q^* < q_{EME}$) implica que:}
\begin{opts}
    \item Las empresas obtienen beneficios económicos extraordinarios positivos de manera persistente gracias a las barreras de entrada.
    \item El precio de equilibrio es estrictamente inferior al costo marginal, generando pérdidas operativas subsidiadas por el Estado.
    \item Cada empresa produce en el tramo decreciente de su curva de costo medio, por debajo de la escala que minimiza el costo medio.
    \item El número de variedades producidas en el mercado es necesariamente menor al que existiría bajo monopolio puro.
\end{opts}

\q{En el modelo de competencia monopolística de Dixit-Stiglitz con preferencias de elasticidad de sustitución constante (CES, $\sigma > 1$) y costos fijos $F$ más costo marginal constante $c$, una expansión del tamaño del mercado (más consumidores):}
\begin{opts}
    \item Aumenta el margen precio-costo de cada firma porque la elasticidad de sustitución $\sigma$ se vuelve menor.
    \item Reduce el número de variedades disponibles en equilibrio al concentrar toda la producción en una sola empresa.
    \item Disminuye el precio hasta el costo marginal $c$ y elimina por completo los costos fijos de producción.
    \item Eleva el número de firmas (variedades) en equilibrio sin alterar el precio ni el margen individual.
\end{opts}

\q{La condición óptima de Dorfman-Steiner para una empresa con poder de mercado que elige simultáneamente precio $P$ y gasto en publicidad $A$ establece que la razón óptima publicidad-ventas ($A / PQ$) es igual a:}
\begin{opts}
    \item El cociente entre la elasticidad-publicidad $\varepsilon_A$ y el valor absoluto de la elasticidad-precio $|\varepsilon_P|$.
    \item El producto de la elasticidad-publicidad de la demanda por el índice de concentración HHI del mercado.
    \item El cociente entre el valor absoluto de la elasticidad-precio $|\varepsilon_P|$ y la elasticidad-publicidad $\varepsilon_A$.
    \item La suma de la elasticidad-precio de la demanda y la elasticidad-publicidad de la demanda de la empresa.
\end{opts}

\q{En el modelo espacial lineal de Hotelling con consumidores distribuidos uniformemente en $[0,1]$ y precios exógenamente fijos e iguales, dos empresas que eligen simultáneamente su ubicación en equilibrio de Nash:}
\begin{opts}
    \item Se ubican en los extremos opuestos ($x_1 = 0$ y $x_2 = 1$) para minimizar los costos totales de transporte social.
    \item Se ubican ambas en el centro del segmento ($x_1 = x_2 = 1/2$), generando diferenciación mínima.
    \item Se ubican en los cuartiles ($x_1 = 1/4$ y $x_2 = 3/4$), logrando exactamente la asignación socialmente óptima.
    \item No alcanzan ningún Equilibrio de Nash en estrategias puras porque siempre prefieren alejarse mutuamente.
\end{opts}

\q{En la extensión de d'Aspremont, Gabszewicz y Thisse (1979) del modelo de Hotelling con costos de transporte cuadráticos $t d^2$ y competencia en precios en la segunda etapa, las empresas eligen en la primera etapa:}
\begin{opts}
    \item Ubicarse juntas en el centro ($x_1 = x_2 = 1/2$) para capturar a todos los consumidores con precios de monopolio.
    \item Ubicarse en los cuartiles ($1/4$ y $3/4$) porque el efecto demanda domina exactamente al efecto estratégico en precios.
    \item Ubicarse en los extremos ($x_1 = 0, x_2 = 1$), logrando diferenciación máxima para suavizar la competencia.
    \item Fijar precios iguales al costo marginal independientemente de la distancia que las separe en el segmento.
\end{opts}

\q{Al comparar el número de empresas que ingresan con libre entrada ($n^*$) frente al número socialmente óptimo ($n^{soc}$) en el modelo circular de Salop con consumidores uniformes y costo fijo de entrada $F$:}
\begin{opts}
    \item El mercado con libre entrada siempre produce insuficientes empresas ($n^* = \frac{1}{2} n^{soc}$) debido a la existencia de costos fijos.
    \item El equilibrio con libre entrada coincide exactamente con el óptimo social ($n^* = n^{soc}$) porque los beneficios son nulos.
    \item El número óptimo social de firmas es infinito porque los costos de transporte tienden a cero sin costos fijos sociales.
    \item La libre entrada produce exceso de firmas ($n^* = 2 n^{soc}$) por la externalidad de robo de negocio.
\end{opts}

\q{En el modelo de diferenciación vertical (calidad) de Shaked y Sutton, cuando los consumidores coinciden en el ordenamiento de calidades pero difieren en su disposición a pagar por calidad:}
\begin{opts}
    \item Las firmas eligen calidades distintas para suavizar precios, y su número puede quedar acotado (finitud).
    \item Todas las empresas producen exactamente la misma calidad máxima para evitar quedarse sin clientes en equilibrio.
    \item La empresa de menor calidad siempre obtiene mayores beneficios que la empresa de alta calidad al tener menores costos.
    \item El equilibrio de mercado converge automáticamente a competencia perfecta con $P = CMg$ en todos los niveles de calidad.
\end{opts}

\q{En el modelo de empresa dominante con franja competitiva (Forchheimer), la elasticidad-precio de la demanda residual que enfrenta la empresa dominante es mayor en valor absoluto cuando:}
\begin{opts}
    \item La elasticidad de la demanda total del mercado es menor, la oferta de la franja es inelástica y la cuota del líder es alta.
    \item La oferta de la franja competitiva es más elástica, la demanda de mercado es más elástica y la cuota del líder es menor.
    \item La cuota de mercado del líder se aproxima al 100\% y existen barreras legales que prohíben la expansión de la franja.
    \item El costo marginal de la empresa dominante es estrictamente creciente e igual al precio fijado por la franja.
\end{opts}

\q{Según la teoría de los mercados perfectamente desafiables (contestable markets) de Baumol, Panzar y Willig, la condición esencial que disciplina incluso a un monopolista único para que fije $P = CMe$ con beneficio nulo es:}
\begin{opts}
    \item La existencia de al menos cinco competidores activos con idéntica estructura de costos en el mercado relevante.
    \item La fijación de un precio tope regulatorio (price cap) por parte del organismo regulador sectorial.
    \item La ausencia de costos hundidos, que permite entrada y salida 'golpea y corre' (hit-and-run).
    \item Que la función de costos presente deseconomías de escala estrictas para cualquier volumen positivo de producción.
\end{opts}

\q{En un monopolio natural multiproducto sin barreras legales de entrada, puede ocurrir que la configuración monopólica sea 'no sostenible' cuando:}
\begin{opts}
    \item Las demandas de los distintos servicios son completamente independientes y presentan elasticidades idénticas.
    \item La función de costos es estrictamente subaditiva y no existen costos hundidos en ningún segmento de la red.
    \item El regulador establece una tarifa igual al costo marginal y financia el déficit mediante transferencias de suma alzada.
    \item Los subsidios cruzados permiten a un entrante practicar 'desnatado' (cream-skimming) en el segmento rentable.
\end{opts}

\q{En el análisis de estructura de mercado, la principal ventaja analítica del Índice de Herfindahl-Hirschman ($HHI = \sum s_i^2$) frente a la razón de concentración de las $k$ mayores firmas ($CR_k$) consiste en que:}
\begin{opts}
    \item El $HHI$ usa todas las firmas y pondera cuadráticamente las cuotas, reflejando asimetrías de tamaño.
    \item El $HHI$ siempre adopta valores estrictamente mayores que $CR_4$ independientemente de la distribución de cuotas.
    \item El $HHI$ requiere únicamente conocer la participación de las dos empresas más grandes del mercado relevante.
    \item El $HHI$ mide directamente las barreras legales de entrada y el grado de diferenciación vertical de los productos.
\end{opts}

\q{En la definición del mercado relevante para el control de fusiones y defensa de la competencia, la prueba del monopolista hipotético (test SSNIP) evalúa si:}
\begin{opts}
    \item Las empresas que desean fusionarse poseen una cuota conjunta superior al 50\% de las ventas nacionales del sector.
    \item Un monopolista hipotético podría elevar el precio en 5\%-10\% de modo no transitorio y rentable.
    \item El costo marginal de la empresa fusionada caería en al menos un 5\% gracias a economías de escala administrativas.
    \item Los consumidores dejarían de comprar por completo el bien ante cualquier incremento de precio superior a la inflación.
\end{opts}

\q{En el sector saneamiento (por ejemplo, concesiones de plantas de tratamiento de aguas residuales PTAR o desalinizadoras), el mecanismo de competencia por el mercado propuesto por Demsetz (1968) consiste en:}
\begin{opts}
    \item Permitir que múltiples operadores tiendan redes paralelas de tuberías en la misma ciudad para competir en precios.
    \item Eliminar toda supervisión regulatoria ex post confiando únicamente en la amenaza de entrada 'golpea y corre'.
    \item Subastar la franquicia exclusiva del monopolio natural al postor que ofrezca la menor tarifa al usuario.
    \item Dividir la empresa prestadora en firmas pequeñas que compitan a la Cournot en la captación de agua cruda.
\end{opts}

\q{Para mitigar el riesgo de renegociaciones oportunistas y de demanda incierta en concesiones de infraestructura (como proyectos de saneamiento APP), la subasta por Menor Valor Presente de los Ingresos (LPVR de Engel-Fischer-Galetovic) se caracteriza porque:}
\begin{opts}
    \item Fija un plazo contractual rígido e inamovible de 30 años y adjudica la obra al postor con mayor capital social.
    \item Garantiza al concesionario una tasa de retorno fija anual sobre sus costos históricos declarados (regulación costo-plus).
    \item Elimina la necesidad de regular la calidad del servicio porque el concesionario compite diariamente en el mercado spot.
    \item Adjudica al menor valor presente de ingresos, con plazo flexible hasta alcanzar dicho monto recaudado.
\end{opts}

\q{En las licitaciones públicas de obras de agua potable y alcantarillado, la práctica anticompetitiva de 'ofertas de resguardo' (cover bidding) consiste en que:}
\begin{opts}
    \item Un postor presenta una propuesta artificialmente alta para simular competencia y dejar ganar a un cómplice.
    \item El organismo convocante fija un valor referencial secreto por debajo del costo marginal para reducir el gasto público.
    \item Dos empresas pequeñas forman un consorcio legítimo y transparente para alcanzar la capacidad técnica mínima exigida en las bases.
    \item El ganador de la licitación subcontrata el 100\% de la obra al Estado a través de un convenio interinstitucional.
\end{opts}

\q{La competencia por comparación (yardstick competition de Shleifer, 1985), fundamento del benchmarking regulatorio que SUNASS aplica entre Empresas Prestadoras de Servicios de Saneamiento (EPS), logra incentivar eficiencia porque:}
\begin{opts}
    \item Permite que las EPS compitan directamente por los mismos usuarios urbanos compartiendo la misma red de distribución.
    \item Desvincula la tarifa reconocida a una EPS de sus propios costos ineficientes, vinculándola al desempeño de empresas comparables.
    \item Asegura que todas las EPS del país cobren exactamente la misma tarifa nominal por metro cúbico sin considerar diferencias geográficas.
    \item Sustituye los estudios tarifarios quinquenales por subastas diarias de derechos de uso de agua en la bolsa de valores.
\end{opts}

\q{En industrias de red con integración vertical, la doctrina de las 'facilidades esenciales' (essential facilities) establece que un propietario verticalmente integrado puede restringir la competencia en el segmento aguas abajo si:}
\begin{opts}
    \item Reduce sus costos operativos en el segmento competitivo mediante innovaciones tecnológicas abiertas al mercado.
    \item Cobra un precio de acceso igual al costo marginal de largo plazo a todos los operadores sin discriminación alguna.
    \item Niega el acceso o aplica tarifas discriminatorias sobre una infraestructura no replicable económicamente.
    \item Permite la desintegración contable y estructural entre la infraestructura de red y la comercialización minorista.
\end{opts}

\q{En el modelo de discriminación de precios en oligopolio de Thisse y Vives (1988), a diferencia del monopolio donde discriminar siempre eleva beneficios, permitir que dos oligopolistas discriminen precios según la ubicación del cliente puede:}
\begin{opts}
    \item Aumentar simultáneamente los precios pagados por todos los consumidores y elevar los beneficios de ambas empresas.
    \item Impedir la existencia de cualquier equilibrio en estrategias puras incluso cuando los costos de transporte son cuadráticos.
    \item Eliminar la diferenciación espacial convirtiendo el mercado en un monopolio natural regulado por el Estado.
    \item Intensificar la competencia en todos los segmentos y reducir los beneficios frente al precio uniforme.
\end{opts}

\q{En mercados oligopólicos con costos de cambio elevados para el consumidor (lock-in), el patrón dinámico de precios típico identificado por Klemperer ('bargain-then-rip-off') consiste en:}
\begin{opts}
    \item Fijar precios bajos iniciales para captar clientes y luego cobrar precios altos a los usuarios cautivos.
    \item Mantener precios exactamente iguales al costo marginal en todos los períodos para evitar la entrada de competidores.
    \item Cobrar precios monopólicos elevados a los clientes nuevos y otorgar descuentos agresivos únicamente a los clientes antiguos.
    \item Fijar precios crecientes únicamente cuando el costo de cambio del consumidor se reduce a cero.
\end{opts}

\q{En un mercado bilateral (two-sided market) con externalidades de red indirectas entre dos grupos de usuarios (por ejemplo, plataformas digitales o tarjetas de pago), la estructura óptima de precios de una plataforma suele implicar:}
\begin{opts}
    \item Cobrar exactamente el mismo precio igual al costo marginal en ambos lados de la plataforma para maximizar el volumen.
    \item Subsidiar al lado más elástico o que genera mayores efectos de red hacia el otro lado de la plataforma.
    \item Fijar precios superiores al costo marginal únicamente cuando la elasticidad de demanda en ambos lados es infinita.
    \item Prohibir el multi-homing en ambos lados para replicar las condiciones de un mercado perfectamente competitivo.
\end{opts}

\q{Al comparar el incentivo a introducir una innovación reductora de costos de un monopolista frente a una empresa en competencia (efecto reemplazo de Arrow, 1962), se concluye que:}
\begin{opts}
    \item Ambos tienen exactamente el mismo incentivo a innovar cuando la innovación es no drástica y patentada perpetuamente.
    \item El monopolista siempre innova más que una firma competitiva porque carece de restricciones de financiamiento externo.
    \item El monopolista innova menos porque la nueva tecnología reemplaza sus rentas monopólicas preexistentes.
    \item Una firma en competencia perfecta nunca tiene incentivo a innovar porque su beneficio tras patentar la innovación sigue siendo nulo.
\end{opts}

\subsection*{Parte II. Aplicaciones prácticas y ejercicios cuantitativos (154 a 180)}

\q{Considere un juego estático $2\times 2$ donde los pagos $(u_1, u_2)$ son: $(A,I)=(4,3)$, $(A,D)=(1,6)$, $(B,I)=(2,5)$ y $(B,D)=(5,2)$. En el Equilibrio de Nash en estrategias mixtas, la probabilidad $p^*$ con la que el Jugador 1 elige $A$ y la probabilidad $q^*$ con la que el Jugador 2 elige $I$ son, respectivamente:}
\begin{opts}
    \item $p^* = 3/4$ y $q^* = 1/3$
    \item $p^* = 2/3$ y $q^* = 1/2$
    \item $p^* = 1/3$ y $q^* = 3/4$
    \item $p^* = 1/2$ y $q^* = 2/3$
\end{opts}

\q{En un juego secuencial con información perfecta, un entrante (Jugador 1) elige entre No Entrar ($N$), Entrada Estándar ($E$) o Entrada Tecnológica ($T$). Si elige $N$, los pagos $(u_1, u_2)$ son $(2, 10)$. Si elige $E$, el incumbente (Jugador 2) elige Acomodar con pagos $(5, 4)$ o Luchar con pagos $(-3, 6)$. Si elige $T$, el incumbente elige Acomodar con pagos $(4, 5)$ o Luchar con pagos $(1, 2)$. El resultado del ENPS es:}
\begin{opts}
    \item Entrada Tecnológica y Acomodar, con pagos $(4, 5)$.
    \item No Entrar al mercado, con pagos finales $(2, 10)$.
    \item Entrada Estándar y Acomodar, con pagos finales $(5, 4)$.
    \item Entrada Tecnológica y Luchar, con pagos finales $(1, 2)$.
\end{opts}

\q{En un oligopolio repetido infinitamente, cada empresa obtiene un beneficio por período de $50$ si coopera en el cártel, $120$ en el período que desvía unilateralmente y $20$ en el equilibrio competitivo de castigo permanente (Grim Trigger). El factor de descuento mínimo $\delta^*$ que sostiene la colusión es:}
\begin{opts}
    \item $\delta^* = 0.50$
    \item $\delta^* = 0.70$
    \item $\delta^* = 0.60$
    \item $\delta^* = 0.80$
\end{opts}

\q{Dos empresas idénticas compiten a la Cournot enfrentando la demanda inversa $P = 120 - 2Q$ (donde $Q = q_1 + q_2$) y un costo marginal constante $CMg = 24$ sin costos fijos. El precio de equilibrio $P^*$ y el beneficio individual $\pi_i^*$ de cada empresa son:}
\begin{opts}
    \item $P^* = 48$ y $\pi_i^* = 384$
    \item $P^* = 64$ y $\pi_i^* = 640$
    \item $P^* = 56$ y $\pi_i^* = 512$
    \item $P^* = 72$ y $\pi_i^* = 576$
\end{opts}

\q{Un mercado oligopólico de Cournot está integrado por $N = 4$ empresas simétricas con costo marginal constante $c = 20$ que enfrentan la demanda inversa $P = 100 - Q$. La cantidad total de equilibrio $Q^*$ y el precio de mercado $P^*$ son:}
\begin{opts}
    \item $Q^* = 80$ y $P^* = 20$
    \item $Q^* = 60$ y $P^* = 40$
    \item $Q^* = 56$ y $P^* = 44$
    \item $Q^* = 64$ y $P^* = 36$
\end{opts}

\q{En un duopolio de Cournot con demanda inversa $P = 90 - Q$, la empresa 1 tiene costo marginal constante $c_1 = 12$ y la empresa 2 tiene $c_2 = 24$. Las cantidades de equilibrio $(q_1^*, q_2^*)$ y el precio $P^*$ son:}
\begin{opts}
    \item $q_1^* = 30$, $q_2^* = 18$ y $P^* = 42$
    \item $q_1^* = 32$, $q_2^* = 16$ y $P^* = 42$
    \item $q_1^* = 28$, $q_2^* = 20$ y $P^* = 42$
    \item $q_1^* = 24$, $q_2^* = 24$ y $P^* = 42$
\end{opts}

\q{Dos empresas simétricas compiten a la Cournot con demanda inversa $P = 90 - (q_1 + q_2)$ y función de costos totales $CT_i(q_i) = 10q_i + 0.5q_i^2$. En equilibrio, la producción individual $q^*$ y el precio $P^*$ son:}
\begin{opts}
    \item $q^* = 16$ y $P^* = 58$
    \item $q^* = 20$ y $P^* = 50$
    \item $q^* = 25$ y $P^* = 40$
    \item $q^* = 30$ y $P^* = 30$
\end{opts}

\q{En un duopolio de Cournot simétrico con demanda $P = 80 - Q$ y costo marginal $c = 20$, el gobierno impone un impuesto específico de $t = 9$ soles por unidad producida. El incremento en el precio de equilibrio ($\Delta P$) es:}
\begin{opts}
    \item $\Delta P = 4.50$ soles
    \item $\Delta P = 7.50$ soles
    \item $\Delta P = 6.00$ soles
    \item $\Delta P = 9.00$ soles
\end{opts}

\q{Una industria de Cournot consta de tres empresas cuyas participaciones de mercado son $50\%$, $30\%$ y $20\%$. Si la elasticidad-precio de la demanda de mercado en equilibrio es $|\varepsilon_d| = 2$, el índice $HHI$ (en escala $0$ a $10\,000$) y el Índice de Lerner promedio ponderado $\bar{L}$ son:}
\begin{opts}
    \item $HHI = 4200$ y $\bar{L} = 0.21$
    \item $HHI = 3800$ y $\bar{L} = 0.38$
    \item $HHI = 3400$ y $\bar{L} = 0.17$
    \item $HHI = 3800$ y $\bar{L} = 0.19$
\end{opts}

\q{En un duopolio de Stackelberg con demanda inversa $P = 100 - 2Q$ y costo marginal constante $c = 20$ para ambas empresas, la cantidad del líder $q_L$, la cantidad del seguidor $q_S$ y el beneficio del líder $\pi_L$ son:}
\begin{opts}
    \item $q_L = 20$, $q_S = 10$ y $\pi_L = 400$
    \item $q_L = 15$, $q_S = 15$ y $\pi_L = 300$
    \item $q_L = 20$, $q_S = 15$ y $\pi_L = 200$
    \item $q_L = 25$, $q_S = 10$ y $\pi_L = 250$
\end{opts}

\q{Un mercado con demanda $P = 100 - Q$ opera como un duopolio de Stackelberg donde la empresa líder tiene costo marginal $c_1 = 16$ y la seguidora tiene $c_2 = 24$. Las cantidades $(q_1^S, q_2^S)$ y el precio $P^S$ en equilibrio son:}
\begin{opts}
    \item $q_1^S = 40$, $q_2^S = 18$ y $P^S = 42$
    \item $q_1^S = 46$, $q_2^S = 15$ y $P^S = 39$
    \item $q_1^S = 44$, $q_2^S = 16$ y $P^S = 40$
    \item $q_1^S = 42$, $q_2^S = 17$ y $P^S = 41$
\end{opts}

\q{Dada la demanda de mercado $P = 120 - 2Q$ y un costo marginal constante $c = 24$ para dos empresas que compiten a la Cournot, la Pérdida de Eficiencia Social ($PES$) respecto al óptimo competitivo es:}
\begin{opts}
    \item $PES = 144$
    \item $PES = 384$
    \item $PES = 256$
    \item $PES = 576$
\end{opts}

\q{Dos empresas venden un bien homogéneo y compiten en precios (Bertrand) sin restricciones de capacidad frente a la demanda $P = 100 - Q$. Si los costos marginales son $c_1 = 20$ y $c_2 = 30$, el beneficio de equilibrio (límite) de la empresa 1 es:}
\begin{opts}
    \item $\pi_1 = 1600$
    \item $\pi_1 = 800$
    \item $\pi_1 = 900$
    \item $\pi_1 = 700$
\end{opts}

\q{Dos empresas producen bienes diferenciados y compiten en precios (Bertrand) con demandas $q_1 = 60 - 2p_1 + p_2$ y $q_2 = 60 - 2p_2 + p_1$, y costo marginal constante $c = 12$. En el Equilibrio de Nash simétrico, el precio $p^*$ y la cantidad $q^*$ de cada firma son:}
\begin{opts}
    \item $p^* = 28$ y $q^* = 32$
    \item $p^* = 24$ y $q^* = 36$
    \item $p^* = 30$ y $q^* = 30$
    \item $p^* = 32$ y $q^* = 28$
\end{opts}

\q{Dos empresas compiten en cantidades (Cournot) ofreciendo bienes diferenciados cuyas demandas inversas son $p_1 = 100 - q_1 - 0.5q_2$ y $p_2 = 100 - q_2 - 0.5q_1$, con costo marginal $c = 25$. La cantidad $q^*$ y el precio $p^*$ de equilibrio de cada empresa son:}
\begin{opts}
    \item $q^* = 25$ y $p^* = 62.5$
    \item $q^* = 30$ y $p^* = 55.0$
    \item $q^* = 35$ y $p^* = 47.5$
    \item $q^* = 40$ y $p^* = 40.0$
\end{opts}

\q{Dos empresas forman un cártel perfecto frente a la demanda $P = 100 - (q_1 + q_2)$, con funciones de costos $CT_1 = q_1^2$ y $CT_2 = 2q_2^2$. Para maximizar el beneficio conjunto, las cuotas $(q_1^*, q_2^*)$ y el precio $P^*$ del cártel deben ser:}
\begin{opts}
    \item $q_1^* = 15$, $q_2^* = 15$ y $P^* = 70$
    \item $q_1^* = 24$, $q_2^* = 12$ y $P^* = 64$
    \item $q_1^* = 20$, $q_2^* = 10$ y $P^* = 70$
    \item $q_1^* = 10$, $q_2^* = 20$ y $P^* = 70$
\end{opts}

\q{En un duopolio de Cournot simétrico con demanda lineal $P = 100 - Q$ y costo marginal $c = 10$ repetido infinitamente, las firmas intentan coludirse produciendo la mitad de la cantidad monopólica cada una y castigan cualquier desvío volviendo permanentemente a Cournot (Grim Trigger). El umbral mínimo $\delta^*$ es:}
\begin{opts}
    \item $\delta^* = 1/2$
    \item $\delta^* = 5/9$
    \item $\delta^* = 8/15$
    \item $\delta^* = 9/17$
\end{opts}

\q{Una empresa en competencia monopolística tiene costos totales $CT(q) = 100 + 20q + q^2$ y en el largo plazo enfrenta la demanda individual $P = 60 - 3q$. La cantidad $q^*$, el precio $P^*$ y el beneficio económico $\pi^*$ de equilibrio son:}
\begin{opts}
    \item $q^* = 5$, $P^* = 45$ y $\pi^* = 0$
    \item $q^* = 5$, $P^* = 45$ y $\pi^* = 15$
    \item $q^* = 6$, $P^* = 42$ y $\pi^* = 0$
    \item $q^* = 10$, $P^* = 30$ y $\pi^* = 0$
\end{opts}

\q{En competencia monopolística, una firma representativa presenta costos $CT(q) = 36 + 10q + q^2$ y enfrenta en el equilibrio de largo plazo la demanda $P = 28 - 1.25q$. El exceso de capacidad ($q_{EME} - q_{LP}^*$) de esta empresa es:}
\begin{opts}
    \item $1.5$ unidades
    \item $2.0$ unidades
    \item $2.5$ unidades
    \item $3.0$ unidades
\end{opts}

\q{Dos empresas ubicadas en los extremos de una ciudad lineal de longitud $1$ (en $x_1 = 0$ y $x_2 = 1$) tienen costo marginal $c = 6$ y sirven a una masa unitaria de consumidores distribuidos uniformemente con costo de transporte $t = 12$ por unidad de distancia. El precio de equilibrio $p^*$ y el beneficio individual $\pi^*$ son:}
\begin{opts}
    \item $p^* = 18$ y $\pi^* = 12$
    \item $p^* = 12$ y $\pi^* = 3$
    \item $p^* = 18$ y $\pi^* = 6$
    \item $p^* = 24$ y $\pi^* = 9$
\end{opts}

\q{En una ciudad circular de perímetro $1$ (modelo de Salop) con masa unitaria de consumidores, el costo de transporte lineal es $t = 16$, el costo marginal es $c = 4$ y el costo fijo de entrada es $F = 1$. El número de empresas en libre entrada $n^*$ y el precio de equilibrio $p^*$ son:}
\begin{opts}
    \item $n^* = 2$ y $p^* = 12$
    \item $n^* = 4$ y $p^* = 10$
    \item $n^* = 8$ y $p^* = 6$
    \item $n^* = 4$ y $p^* = 8$
\end{opts}

\q{En un mercado con demanda total $D(P) = 120 - 2P$, existe una empresa dominante con $CMg_d = 15$ y una franja competitiva cuya oferta agregada es $S_f(P) = P - 15$. El precio de equilibrio $P^*$, la cantidad de la dominante $q_d$ y la cantidad de la franja $q_f$ son:}
\begin{opts}
    \item $P^* = 30$, $q_d = 45$ y $q_f = 15$
    \item $P^* = 35$, $q_d = 30$ y $q_f = 20$
    \item $P^* = 25$, $q_d = 60$ y $q_f = 10$
    \item $P^* = 30$, $q_d = 35$ y $q_f = 25$
\end{opts}

\q{Un mercado está compuesto por cuatro empresas con participaciones de $40\%$, $25\%$, $20\%$ y $15\%$. Si las dos empresas más pequeñas deciden fusionarse (manteniendo constantes las participaciones), el incremento $\Delta HHI$ y el índice $HHI_{\text{post}}$ resultante son:}
\begin{opts}
    \item $\Delta HHI = 300$ y $HHI_{\text{post}} = 3150$
    \item $\Delta HHI = 600$ y $HHI_{\text{post}} = 3450$
    \item $\Delta HHI = 600$ y $HHI_{\text{post}} = 3250$
    \item $\Delta HHI = 700$ y $HHI_{\text{post}} = 3550$
\end{opts}

\q{En un mercado de Cournot con $3$ empresas simétricas, demanda $P = 120 - Q$ y costo marginal $c = 0$, dos de las empresas se fusionan sin generar sinergias de costos. Comparando el beneficio de la empresa fusionada frente a la suma de beneficios previos de ambas, la variación es:}
\begin{opts}
    \item Una ganancia neta de $+200$ unidades monetarias.
    \item Una ganancia neta de $+400$ unidades monetarias.
    \item Una pérdida neta de $-200$ unidades monetarias.
    \item Una pérdida neta de $-400$ unidades monetarias.
\end{opts}

\q{Un fabricante monopolista con costo marginal $c = 20$ vende mediante un precio lineal mayorista $w$ a un distribuidor minorista monopolista (sin costos adicionales) que enfrenta la demanda final $P = 100 - Q$. El precio final con doble marginalización ($P_{DM}$) y el precio bajo integración vertical ($P_{IV}$) son:}
\begin{opts}
    \item $P_{DM} = 70$ y $P_{IV} = 50$
    \item $P_{DM} = 80$ y $P_{IV} = 50$
    \item $P_{DM} = 75$ y $P_{IV} = 60$
    \item $P_{DM} = 80$ y $P_{IV} = 60$
\end{opts}

\q{Una empresa oligopólica con diferenciación de producto estima que la elasticidad-precio de su demanda es $|\varepsilon_P| = 2.5$ y la elasticidad-publicidad de su demanda es $\varepsilon_A = 0.25$. Según la condición de Dorfman-Steiner, la proporción óptima de gasto en publicidad respecto a sus ingresos por ventas ($A / PQ$) es:}
\begin{opts}
    \item $10.0\%$ de las ventas
    \item $12.5\%$ de las ventas
    \item $20.0\%$ de las ventas
    \item $25.0\%$ de las ventas
\end{opts}

\q{En una subasta de sobre cerrado de primer precio participan $N = 3$ postores neutrales al riesgo cuyas valoraciones privadas independientes se distribuyen uniformemente en $[0, 100]$. Si un postor posee una valoración privada de $v_i = 75$, su puja óptima en el Equilibrio Bayesiano de Nash simétrico es:}
\begin{opts}
    \item $b^*(75) = 25.0$
    \item $b^*(75) = 50.0$
    \item $b^*(75) = 37.5$
    \item $b^*(75) = 75.0$
\end{opts}
